\subsubsection{APIs externas}

La arquitectura de \textit{APIs externas} en el objetivo final de este proyecto consiste en un único programa en el lenguaje de programación \textbf{Python}, 
el cual será capaz de extraer información de forma precisa y paralela de múltiples bases de datos accesibles mediante APIs con una interfaz sencilla e intuitiva, 
con el fin de permitir seleccionar los datos deseados de forma filtrada, sin necesidad de filtrar los datos almacenados, 
reduciendo la carga en nuestros servidores a cambio de un incremento en el coste de comunicaciones y del coste en los servidores de la API. 

Además, se planea automatizar mediante funciones más avanzadas y con asistencia de IA las llamadas de API, para permitir llamadas más precisas y sencillas, pero increíblemente eficaces.

Los datos de los objetos extraídos son analizados, y si se consideran relevantes se puede adaptar fácilmente la estructura de los objetos de la base de datos, 
gracias a emplear una estructura flexible y dinámica tanto en la base de datos como en la extracción de los mismos. 
Mediante este análisis se puede además detectar datos incorrectos o inválidos, para mantener una base de datos de calidad.

Este método de obtención de datos se realizará cuando se requiera, no de forma periódica, debido al gran rango de tamaño de datos que se descargan, acorde a la petición realizada . 
Para permitir actualizaciones frecuentes, se ha empleado una arquitectura muy eficiente, capaz de procesar cantidades masivas de datos en segundos, aunque limitada por la velocidad de comunicación y de procesamiento del server API.
